\documentclass[11pt]{letter}
\usepackage[margin=1in]{geometry}
\usepackage{hyperref}
\signature{Julian Juan}
\address{Julian Juan\\Independent researcher/scholar\\\href{mailto:juliansjuan08@gmail.com}{juliansjuan08@gmail.com}}
\begin{document}

\begin{letter}{Editors\\Eurosurveillance}

\sloppy
\opening{Dear Editors,}

Please consider the manuscript entitled ``Country-level reported hantavirus incidence in the EU/EEA, 2019--2023: an open benchmark with public covariates and probabilistic evaluation'' for publication as a regular surveillance/research article in \emph{Eurosurveillance}.

The manuscript presents a reproducible ECDC-only EU/EEA country-year benchmark for reported hantavirus infection incidence from 2019 through 2023. The work links ECDC annual surveillance labels to public World Bank, FAOSTAT, TerraClimate, and Natural Earth sources, evaluates retrospective one-year-ahead probabilistic predictions, and reports calibration and sharpness alongside weighted interval score.

The framing is intentionally conservative. The manuscript does not present a live outbreak tracker, does not claim prospective forecasting, does not make county-level or individual-risk claims, and does not claim causal climate or vegetation effects. Its main contribution is a transparent surveillance evaluation showing that simple empirical baselines remain difficult to beat under sparse annual country-level public data. This negative and mixed result is important because it helps prevent false precision and over-sensationalized interpretation of rare but severe zoonotic infections.

All data used are aggregated and publicly available. No individual-level patient data were accessed. No external funding was received, and I declare no competing interests.

\closing{Sincerely,}

\end{letter}
\end{document}
